% Generated draft: requires mathematical and source review.
\documentclass[11pt]{article}
\usepackage[margin=25mm]{geometry}
\usepackage{fontspec}
\setmainfont{Latin Modern Roman}
\usepackage{amsmath,amssymb,mathtools,unicode-math}
\setmathfont{Latin Modern Math}
\usepackage{enumitem,needspace,xurl,hyperref}
\setlength{\parindent}{0pt}
\setlength{\parskip}{5pt}
\setlength{\emergencystretch}{3em}
\newcommand{\N}{\mathbb{N}}
\newcommand{\Z}{\mathbb{Z}}
\newcommand{\Q}{\mathbb{Q}}
\newcommand{\R}{\mathbb{R}}
\newcommand{\C}{\mathbb{C}}
\newcommand{\F}{\mathbb{F}}
\newcommand{\A}{\mathbb{A}}
\newcommand{\PP}{\mathbb{P}}
\newcommand{\E}{\mathbb{E}}
\newcommand{\eps}{\varepsilon}
\newcommand{\dd}{\,\mathrm{d}}
\begin{document}
\section*{UnsolvedMath ID: 801083}
\textbf{Source problem code:} catalogue entry
\subsection{Definitions and mathematical statement}
The supplied record does not define the underlying dynamical system, the potential $\phi$, or the measure $\mu_\phi$. Consequently, the following expresses the requested conclusion, with these objects left subject to their missing source hypotheses.

Write $f\colon X\to X$ for the intended underlying map. For a real-valued observable $g$ for which the integral exists, set
\[
S_n(g)(x)=\sum_{j=0}^{n-1}g(f^j(x)),\qquad
\langle\mu_\phi,g\rangle=\int_X g\,\mathrm d\mu_\phi,
\qquad
Z_n(g)=\frac{S_n(g)-n\langle\mu_\phi,g\rangle}{\sqrt n}.
\]
These formulas interpret the notation in the record; the map, its domain, the probability measure, and the integrability hypotheses must be recovered from the source context. Relative to the intended metric $d$ on $X$, a real-valued observable is H\"older continuous if there are constants $C\geq0$ and $\alpha>0$ such that
\[
|g(x)-g(y)|\leq C\,d(x,y)^\alpha\qquad(x,y\in X).
\]
The intended metric, observable domain, and any boundedness or local versus global H\"older convention are not supplied.

\textbf{Requested assertion, with unresolved source qualifications.} For every intended H\"older continuous observable $g$, there exists $\sigma=\sigma(g)\geq0$ such that, if $\sigma^2>0$, then for every interval $I\subseteq\R$,
\[
\lim_{n\to\infty}\mu_\phi\bigl(\{x\in X:Z_n(g)(x)\in I\}\bigr)
=\frac{1}{\sqrt{2\pi\sigma^2}}\int_I
\exp\!\left(-\frac{t^2}{2\sigma^2}\right)\,\mathrm dt.
\]
If $\sigma^2=0$, the displayed source conclusion specifies only
\[
\lim_{n\to\infty}\mu_\phi\bigl(\{x\in X:Z_n(g)(x)\in(-\delta,\delta)\}\bigr)=1
\qquad(\delta>0).
\]
The positivity condition on $\delta$ makes explicit the neighborhood interpretation of this clause. Although the source introduces its formula with a quantifier over every interval, its zero-variance branch gives no conclusion for other intervals. A standard central limit theorem would express convergence in distribution to $\mathcal N(0,\sigma^2)$, interpreting $\mathcal N(0,0)$ as the point mass at zero; this interpretation does not imply convergence of probabilities for every interval with zero on its boundary. The intended zero-variance formulation therefore requires clarification.
\subsection{Short English statement}
The problem asks whether the intended measure $\mu_\phi$ satisfies a central limit theorem for every real-valued H\"older observable, with variance allowed to vanish. The supplied excerpt omits the dynamical and measure hypotheses and leaves its assertion for arbitrary intervals incomplete in the zero-variance case.
\subsection{Sources}
\begin{itemize}
\item Local UnsolvedMath registry: permanent ID $801083$, exact source problem code \texttt{catalogue entry}, title ``Henon Maps: Open Problems --- Problem 35.'' The record attributes the problem to Berger and Bedford, with proposed year $2023$.
\item Berger, Bedford, \emph{Henon Maps: A List of Open Problems} (2023), Problem 35. Supplied source URL: \url{https://arxiv.org/abs/2312.03907}. The record reports extraction by \texttt{arxiv\_latex} and access on 2026-10-16; these are supplied provenance claims, not independently verified access or source verification.
\item The local record labels the problem \texttt{open}. This is a supplied status label, not an independently verified assessment of the literature.
\end{itemize}
\end{document}
